% ===============================================================================
% ML-06b: Musterlösung - ir + estar (Bewegung und Lage)
% Abgeleitet aus LE-06b (Abschnitt 9: Lösungen)
% ===============================================================================

\documentclass[11pt, a4paper]{article}

\usepackage[utf8]{inputenc}
\usepackage[T1]{fontenc}
\usepackage[ngerman]{babel}
\usepackage{helvet}
\renewcommand{\familydefault}{\sfdefault}

\usepackage[margin=2cm]{geometry}
\usepackage{enumitem}
\usepackage{tabularx}
\usepackage{booktabs}

\usepackage{xcolor}
\usepackage{tcolorbox}

\usepackage{fancyhdr}
\usepackage{lastpage}

% Farben
\definecolor{spanisch}{HTML}{c41e3a}
\definecolor{grau}{HTML}{888888}
\definecolor{hellgrau}{HTML}{f5f5f5}
\definecolor{orange}{HTML}{b35c00}
\definecolor{loesung}{HTML}{c62828}

\newcommand{\fachfarbe}{spanisch}

% Variablen
\newcommand{\thema}{ir + estar (Bewegung und Lage)}
\newcommand{\sequenz}{06}
\newcommand{\block}{b}
\newcommand{\fach}{Spanisch}
\newcommand{\klasse}{7}
\newcommand{\maxpunkte}{20}

% Kopf-/Fußzeile
\pagestyle{fancy}
\fancyhf{}
\renewcommand{\headrulewidth}{0.4pt}
\renewcommand{\footrulewidth}{0.4pt}
\lhead{\fach{} -- Klasse \klasse}
\chead{\textcolor{loesung}{\textbf{ML-\sequenz\block{} (MUSTERLÖSUNG)}}}
\rhead{}
\lfoot{}
\rfoot{Seite \thepage\ von \pageref{LastPage}}

% Befehle
\newcommand{\punktebox}[1]{\hfill\fbox{\small \_/#1}}
\newcommand{\loesung}[1]{\textcolor{loesung}{\textbf{#1}}}
\newcommand{\luecke}[1]{\rule{#1}{0.4pt}}
\newcommand{\es}[1]{\textcolor{\fachfarbe}{\textbf{#1}}}

% Merkkasten
\newtcolorbox{merkkasten}[1][]{
    colback=hellgrau,
    colframe=\fachfarbe,
    boxrule=1pt,
    arc=0pt,
    left=8pt, right=8pt, top=8pt, bottom=8pt,
    fonttitle=\bfseries\color{\fachfarbe},
    title=#1
}

% Hinweis
\newtcolorbox{hinweis}{
    colback=white,
    colframe=orange,
    boxrule=1pt,
    arc=0pt,
    left=5pt, right=5pt, top=5pt, bottom=5pt,
    fonttitle=\bfseries,
    title=Hinweis
}

% Aufgabe
\newenvironment{aufgabe}[1]{%
    \vspace{10pt}
    \noindent\textbf{\textcolor{\fachfarbe}{Aufgabe #1}}
    \vspace{5pt}
    \begin{list}{}{\leftmargin=0pt}
    \item[]
}{%
    \end{list}
}

% ===============================================================================
\begin{document}

% Kopfbereich
\begin{center}
    {\Large\bfseries\textcolor{\fachfarbe}{Arbeitsblatt: \thema}}\\[0.3em]
    {\large\textcolor{loesung}{\textbf{--- MUSTERLÖSUNG ---}}}
\end{center}

\vspace{5pt}

\noindent
\begin{tabularx}{\textwidth}{|X|X|X|}
\hline
\textbf{Name:} \textcolor{loesung}{Musterschüler/in} &
\textbf{Klasse:} \klasse &
\textbf{Datum:} \textcolor{loesung}{XX.XX.XXXX} \\
\hline
\end{tabularx}

\vspace{15pt}

% ===============================================================================
% MERKKASTEN 1: Verb ir (mit Lösungen)
% ===============================================================================

\begin{merkkasten}[Merkkasten 1: Das Verb \textit{ir} (gehen/fahren)]

\textbf{Struktur:} ir + a + Ort

\begin{center}
\begin{tabular}{|l|l|l|}
\hline
\textbf{Person} & \textbf{Pronomen} & \textbf{ir} \\
\hline
1. Sg. & yo & \loesung{voy} \\
\hline
2. Sg. & tú & \loesung{vas} \\
\hline
3. Sg. & él/ella & \loesung{va} \\
\hline
1. Pl. & nosotros & \loesung{vamos} \\
\hline
2. Pl. & vosotros & \loesung{vais} \\
\hline
3. Pl. & ellos/ellas & \loesung{van} \\
\hline
\end{tabular}
\end{center}

\vspace{0.5em}
\textbf{Wichtig:} \es{a + el = al} \hspace{1cm} Beispiel: Voy \textbf{al} cine. (nicht: a el cine)
\end{merkkasten}

\vspace{10pt}

% ===============================================================================
% MERKKASTEN 2: Verb estar (mit Lösungen)
% ===============================================================================

\begin{merkkasten}[Merkkasten 2: Das Verb \textit{estar} (sein/sich befinden)]

\textbf{Verwendung:} Ortsangaben -- Wo befindet sich etwas?

\begin{center}
\begin{tabular}{|l|l|l|}
\hline
\textbf{Person} & \textbf{Pronomen} & \textbf{estar} \\
\hline
1. Sg. & yo & \loesung{estoy} \\
\hline
2. Sg. & tú & \loesung{estás} \\
\hline
3. Sg. & él/ella & \loesung{está} \\
\hline
1. Pl. & nosotros & \loesung{estamos} \\
\hline
2. Pl. & vosotros & \loesung{estáis} \\
\hline
3. Pl. & ellos/ellas & \loesung{están} \\
\hline
\end{tabular}
\end{center}

\vspace{0.5em}
\textbf{Beispiel:} El cine \es{está} en la plaza. = Das Kino ist auf dem Platz.
\end{merkkasten}

\vspace{15pt}

% ===============================================================================
% AUFGABE 1 (mit Lösungen)
% ===============================================================================

\begin{aufgabe}{1}
Setze die richtige Form von \textit{ir} ein. \punktebox{4}
\end{aufgabe}

\begin{enumerate}[label=\alph*)]
    \item Yo \loesung{voy} al cine.
    \item Tú \loesung{vas} a la plaza.
    \item María \loesung{va} al parque.
    \item Nosotros \loesung{vamos} a la biblioteca.
\end{enumerate}

\textit{\textcolor{grau}{Je 1P pro richtige Form}}

\vspace{15pt}

% ===============================================================================
% AUFGABE 2 (mit Lösungen)
% ===============================================================================

\begin{aufgabe}{2}
Ergänze mit \textit{al} oder \textit{a la}. Beachte die Kontraktion! \punktebox{4}
\end{aufgabe}

\begin{enumerate}[label=\alph*)]
    \item Voy \loesung{al} supermercado. (a + el = al)
    \item Vas \loesung{a la} plaza. (la $\rightarrow$ keine Kontraktion)
    \item Va \loesung{al} parque. (a + el = al)
    \item Vamos \loesung{a la} biblioteca. (la $\rightarrow$ keine Kontraktion)
\end{enumerate}

\textit{\textcolor{grau}{Je 1P pro richtige Form. Kontraktion \textit{al} nur bei maskulinen Nomen!}}

\vspace{15pt}

% ===============================================================================
% AUFGABE 3 (mit Lösungen)
% ===============================================================================

\begin{aufgabe}{3}
Setze die richtige Form von \textit{estar} ein und beschreibe den Ort. \punktebox{6}
\end{aufgabe}

\begin{enumerate}[label=\alph*)]
    \item El cine \loesung{está} en la plaza.
    \item La biblioteca \loesung{está} cerca del parque.
    \item ¿Dónde \loesung{estás} (tú)? -- \loesung{Estoy} en casa.
\end{enumerate}

\textit{\textcolor{grau}{Je 1,5P pro Lücke. Achte auf die richtige Person!}}

\vspace{15pt}

% ===============================================================================
% AUFGABE 4 (Beispiellösung)
% ===============================================================================

\begin{aufgabe}{4}
Schreibe einen Dialog. Verwende \textit{ir} und \textit{estar}. (6 Zeilen) \punktebox{6}
\end{aufgabe}

\textbf{Beispiellösung:}

\loesung{A: ¡Hola! ¿Adónde vas?}

\loesung{B: Voy al cine. ¿Y tú?}

\loesung{A: Voy a la plaza. ¿Dónde está el cine?}

\loesung{B: El cine está en la calle Mayor.}

\loesung{A: ¡Gracias! ¡Hasta luego!}

\loesung{B: ¡Adiós!}

\vspace{0.5em}
\textit{\textcolor{grau}{Bewertung: ir-Form (2P), estar-Form (2P), Struktur/Inhalt (1P), Korrektheit (1P)}}

\vspace{20pt}
\hrule
\vspace{10pt}

\begin{flushright}
\textbf{Punkte:} \loesung{\maxpunkte} / \maxpunkte
\end{flushright}

\end{document}
